\scanpage{16}
\mspage{3}
\textcircled{3} Soit
\deleted{$X\supset Y$, $X$ muni de sous-espaces}
\inserted{$\bar T$ un vois. \deleted{tub.} conique absolu de $Y$.}
Soit $\bar T'\deleted{\unsure{_\varepsilon}}\subset\bar T$.
On dit que $\bar T'$ est \uline{sous-espace conique} dans $\bar T$ si on a
\[
(12)\qquad\left\{\begin{array}{ll}
a)&\bar T'\text{ stable par homothéties }\lambda\in I\
\deleted{\unsure{\lambda\bar T'\subset\bar T'}}\\
b)&\text{Pour cette opération, }\bar T'\text{ est un voisinage}\\
&\deleted{\text{tubulaire}}\text{ conique de }Y.
\end{array}\right.
\]
\uline{NB} On aura nécessairement $\dot T'=\bar T'\cap\dot T$,
$\bar T'_\lambda=\bar T'\cap\bar T_\lambda$ ;
\deleted{Si $(\bar T_i)_{i\in I}$ est une famille}
\inserted{la donnée de $\bar T'$ équivaut à}
celle d'un sous-espace fermé $\dot T'$ de $\dot T$
et d'un sous-espace $Y'$ de $Y$, tel que $Y'\supset0.\dot T'$, et
on a
\[
\bar T'=\operatorname{Im}(I\times\dot T'
\deleted{\unsure{_\lambda}}\longrightarrow\bar T)\cup Y'.
\]
\margin{On dit que $\bar T'$ est conique \uline{pur}
si $\bar T'\cap Y=0.\dot T'$,
--- \unsure{ces $\bar T'$ correspondent donc aux fermés}
$\dot T'$ de $\dot T$.}

On dira qu'une structure conique sur
un $\bar T\supset Y$ est \uline{compatible} avec des sous-espaces
\inserted{fermés} $\bar T_i$ de $\bar T$ $(i\in I)$ si ces derniers
sont coniques pour cette structure.

Soit $\bar T$ un voisinage conique de $Y\subset X$
dans $X$, soit $X'$ partie fermée de $X$. On
dit que $\bar T$ est \uline{compatible} avec $X'$ si
$X'\cap\bar T$ est conique dans $\bar T$. On dit qu'il
est \uline{strictement compatible}$^{*}$ avec $\bar T$, si
$\bar T$ est voisinage conique strict, et s'il existe
un voisinage cylindrique de $\dot T$ dans
$X^*=X\setminus\mathring T$, compatible avec
$X'{}^*=X'\cap X^*$,
ou ce qui revient au $=$, s'il existe un
voisinage $E$ de $\dot T$ dans $X^*$, et un iso
\margin{\deleted{\unsure{Définir la notion pour $X'$}}}
\margin{$^{*}$ \deleted{\uline{\unsure{aussi dire}} \uline{\unsure{strict}}}}
\margin{modifier \unsure{sans doute} la définition des
voisinages coniques stricts \unsure{\ldots}}
\margin{15}
